\section{Bài toán và bộ dữ liệu}
\label{sec:problem}

\subsection{Phát biểu bài toán}

Đề tài giải quyết bài toán phân loại âm thanh đa lớp có giám sát. Cho một đoạn
âm thanh ngắn, mô hình cần dự đoán nhãn của đoạn đó trong số mười lớp âm thanh đô
thị đã định trước. Gọi $x$ là đoạn âm thanh đầu vào và $y$ là nhãn tương ứng, mô
hình học một hàm ánh xạ
\[
f(x) \to y, \qquad y \in \{1, 2, \dots, 10\}.
\]
Đầu ra của mô hình là một phân phối xác suất trên mười lớp; nhãn dự đoán là lớp
có xác suất cao nhất. Đây là bài toán \textit{supervised multi-class
classification}: từ tập huấn luyện gồm các cặp $(x_i, y_i)$, mô hình được tối ưu
để giảm sai số dự đoán, đồng thời phải tổng quát hóa tốt sang dữ liệu chưa thấy.

\subsection{Ý nghĩa của bài toán}

Trong đô thị, âm thanh là một kênh thông tin giàu giá trị nhưng thường bị bỏ qua.
Một hệ thống giám sát có thể dựa vào tiếng súng hay còi báo động để phát cảnh
báo; phân tích giao thông có thể tận dụng tiếng còi xe và động cơ; quan trắc môi
trường có thể ước lượng mức ồn và thành phần âm thanh đặc trưng của một khu vực.
Về mặt học thuật, đây là bài toán tiêu biểu để kiểm chứng năng lực của mô hình
học sâu trên dữ liệu phi ảnh, vì nó đòi hỏi kết hợp xử lý tín hiệu số, biểu diễn
đặc trưng và mô hình hóa -- ba mảng kiến thức cốt lõi của môn học.

\subsection{Lý do chọn bộ dữ liệu UrbanSound8K}

UrbanSound8K là một trong những bộ dữ liệu chuẩn cho phân loại âm thanh môi
trường. Đề tài chọn bộ dữ liệu này vì ba lý do thực tế. Dữ liệu đã được gán nhãn
sẵn theo từng lớp, phù hợp trực tiếp với bài toán có giám sát. Bộ dữ liệu cũng
được tác giả gốc chia sẵn thành 10 fold, giúp tổ chức tập huấn luyện, kiểm định
và kiểm thử một cách rõ ràng mà vẫn giữ được tính tách biệt. Cuối cùng, kích
thước dữ liệu vừa phải -- đủ để huấn luyện CNN và ResNet trên một GPU đơn, nhưng
không quá nhỏ để khiến kết quả mất ý nghĩa.

\subsection{Tổng quan bộ dữ liệu}

Bước phân tích dữ liệu khám phá (EDA) cho thấy UrbanSound8K gồm 8732 đoạn âm
thanh, chia đều về mặt tổ chức thành 10 fold và 10 lớp. Thời lượng trung bình của
mỗi đoạn khoảng 3{,}61 giây, nhưng phân bố rất lệch: phần lớn đoạn dài xấp xỉ 4
giây, trong khi một số đoạn chỉ dài 0{,}05 giây. Hai yếu tố không đồng nhất khác
cũng nổi lên rõ. Thứ nhất, tần số lấy mẫu trải dài từ 8\,kHz đến 192\,kHz, trong
đó 44{,}1\,kHz và 48\,kHz chiếm đa số. Thứ hai, dữ liệu trộn lẫn cả tín hiệu mono
và stereo, với phần lớn là stereo.

\begin{figure}[H]
  \centering
  \imgfallback{eda_summary.png}{0.62\textwidth}
  \caption{Tóm tắt kết quả EDA: số mẫu, số lớp, phân bố tần số lấy mẫu, số kênh
           và thống kê thời lượng.}
  \label{fig:eda}
\end{figure}

Mười lớp âm thanh gồm: \texttt{air\_conditioner}, \texttt{car\_horn},
\texttt{children\_playing}, \texttt{dog\_bark}, \texttt{drilling},
\texttt{engine\_idling}, \texttt{gun\_shot}, \texttt{jackhammer}, \texttt{siren}
và \texttt{street\_music}.

\subsubsection{Đặc điểm âm học của các lớp}

Mười lớp này không đồng đều về độ khó. Một nhóm lớp có đặc trưng âm học nổi bật
và dễ tách: \texttt{car\_horn}, \texttt{gun\_shot} và \texttt{siren} đều có cấu
trúc năng lượng tập trung hoặc thay đổi mạnh, đột ngột theo thời gian.
\texttt{dog\_bark} cũng tương đối dễ nhờ các xung năng lượng rõ.

Ngược lại, nhóm âm thanh \textit{cơ khí, mang tính nền} là nguồn nhầm lẫn chính.
\texttt{air\_conditioner}, \texttt{engine\_idling}, \texttt{drilling} và
\texttt{jackhammer} đều là những âm thanh lặp lại, năng lượng phân bố khá đều và
phổ tần số gần nhau. Một con người nghe kỹ vẫn phân biệt được, nhưng trên
spectrogram, ranh giới giữa chúng khá mờ. Quan sát này được nhắc lại nhiều lần ở
phần phân tích kết quả, vì đúng nhóm lớp này quyết định phần lớn sai số của mọi
mô hình. Nó cũng là lý do đề tài không chỉ dùng Accuracy mà bổ sung Macro-F1 và
confusion matrix khi đánh giá.

\subsection{Kết quả phân tích dữ liệu khám phá}

\subsubsection{Tính nhất quán giữa metadata và file âm thanh}

Trước khi xây dựng pipeline, đề tài kiểm tra xem mỗi dòng metadata có ánh xạ đúng
tới một file âm thanh tồn tại trên đĩa hay không. Kết quả cho thấy toàn bộ ánh xạ
đều chính xác, không có file thiếu hay sai đường dẫn. Đây là điều kiện cần để các
bước \texttt{Dataset}/\texttt{DataLoader} về sau không bị sai nhãn.

\subsubsection{Tần số lấy mẫu}

Tần số lấy mẫu ảnh hưởng trực tiếp tới dải tần biểu diễn được và tới số mẫu tín
hiệu theo thời gian. Nếu để nguyên sự chênh lệch từ 8\,kHz đến 192\,kHz, mô hình
có nguy cơ học nhầm \textit{sự khác biệt về định dạng} thay vì đặc trưng âm học
của lớp. Do đó toàn bộ dữ liệu cần được đưa về một tần số chung.

\begin{figure}[H]
  \centering
  \imgfallback{eda_sample_rate.png}{0.7\textwidth}
  \caption{Phân bố tần số lấy mẫu của UrbanSound8K. Phần lớn đoạn ở mức
           44{,}1\,kHz và 48\,kHz.}
  \label{fig:sr}
\end{figure}

\subsubsection{Số kênh}

Vì dữ liệu trộn mono và stereo, đầu vào mô hình sẽ không đồng nhất nếu giữ nguyên
số kênh. Việc gộp về mono vừa chuẩn hóa kích thước đầu vào, vừa giảm chi phí tính
toán mà gần như không làm mất thông tin phân loại.

\begin{figure}[H]
  \centering
  \imgfallback{eda_channels.png}{0.55\textwidth}
  \caption{Phân bố số kênh: đa số đoạn là stereo, một phần nhỏ là mono.}
  \label{fig:ch}
\end{figure}

\subsubsection{Thời lượng}

Độ dài các đoạn không cố định, tuy đa số gần 4 giây. Do CNN và ResNet đòi hỏi đầu
vào có kích thước cố định, cần một chiến lược chuẩn hóa thời lượng. Đề tài chọn
mốc 4{,}0 giây -- vừa bao phủ phần lớn các đoạn mà không phải cắt bỏ nhiều, vừa
là một con số tròn thuận tiện cho việc tính toán kích thước đặc trưng.

\begin{figure}[H]
  \centering
  \imgfallback{eda_duration.png}{0.7\textwidth}
  \caption{Phân bố thời lượng: phần lớn đoạn dồn quanh mốc 4 giây.}
  \label{fig:dur}
\end{figure}

\subsection{Quyết định tiền xử lý rút ra từ EDA}

Kết hợp các quan sát trên, đề tài cố định bốn quyết định tiền xử lý: chuyển toàn
bộ tín hiệu về mono; resample về 22{,}05\,kHz; chuẩn hóa độ dài về 4{,}0 giây
bằng cách pad hoặc cắt; và trích xuất log-Mel spectrogram làm đặc trưng đầu vào.
Bốn quyết định này được mô tả chi tiết ở Mục~\ref{sec:method} và được áp dụng
giống hệt nhau cho mọi mô hình, để khác biệt về kết quả chỉ phản ánh khác biệt về
mô hình chứ không phải về dữ liệu.
